\hypertarget{positive-lower-type-and-quantitative-tap-continuity}{%
\chapter[Growth and tap continuity]{Positive lower type and quantitative tap
continuity}\label{app:growth}\label{positive-lower-type-and-quantitative-tap-continuity}}

This chapter sharpens the entire-growth theorem for nonlinear contraction families. The symmetric definition has the antecedent identified in the reading map.

\hypertarget{the-exact-growth-exponent-has-positive-lower-type}{%
\section{The exact growth exponent has positive lower
type}\label{the-exact-growth-exponent-has-positive-lower-type}}

Let

\[
F'(z)=1+\sum_{j=1}^{J}h_jF(q_jz)F(r_jz),\qquad F(0)=0,
\]

where \(h_j>0\), \(\sum_jh_j=1\), and \(0<q_j,r_j<1\). Zero-weight
channels may instead be allowed if they are removed from the active bank
before defining its growth exponent. Put

\[
\rho=\max_{j:h_j>0}\rho_j,\qquad q_j^{\rho_j}+r_j^{\rho_j}=1.
\]

The inherited coefficient recurrence gives nonnegative coefficients,
entire extension, and \(M(R)=\max_{|z|=R}|F(z)|=F(R)\ge R\) for \(R>0\).

\textbf{Theorem A.1.} There are explicit constants \(0<c_-\le c_+<\infty\)
such that

\[
c_-\le\liminf_{R\to\infty}\frac{\log M(R)}{R^\rho}
\le\limsup_{R\to\infty}\frac{\log M(R)}{R^\rho}\le c_+.
\]

In particular the bank has positive lower type and finite upper type at
its exact order, including when \(0<\rho<1\). Furthermore

\[
\lim_{R\to\infty}\frac{\log\log M(R)}{\log R}=\rho.
\]

This last radial growth-exponent limit does not assert a type limit or
uniform angular growth.

\hypertarget{explicit-lower-constants-and-proof}{%
\subsection{Explicit lower constants and
proof}\label{explicit-lower-constants-and-proof}}

Select any active maximizing channel and abbreviate it to \((h,q,r)\),
so \(q^\rho+r^\rho=1\). Define

\[
L=\frac1h,\qquad
C=\max\left(\frac{qL}{1-q},\frac{rL}{1-r}\right),\qquad
B=L+\frac2{\min(q,r)},\qquad
c_-=\frac{\log2}{(B+C)^\rho}.
\]

We prove the stronger bound

\[
\boxed{\log M(R)\ge c_-(R+C)^\rho\quad\text{for every }R\ge2.}
\]

Write \(H(R)=\log M(R)\). On the seed interval \([2,B]\),
\(H(R)\ge\log R\ge\log2\ge c_-(R+C)^\rho\). For \(R\ge B\), integrate
the positive feedback over the fixed-width interval \([R-L,R]\):

\[
\begin{aligned}
M(R)&\ge h\int_{R-L}^{R}M(qt)M(rt)\,dt\\
&\ge hL M(q(R-L))M(r(R-L))\\
&=M(q(R-L))M(r(R-L)).
\end{aligned}
\]

The two child radii are at least two because \(R\ge B\). Moreover

\[
q(R-L)+C\ge q(R+C),\qquad
r(R-L)+C\ge r(R+C),
\]

because respectively \((1-q)C\ge qL\) and \((1-r)C\ge rL\). If the
comparison bound is known at both child radii, it follows that

\[
\begin{aligned}
H(R)&\ge H(q(R-L))+H(r(R-L))\\
&\ge c_-\{[q(R-L)+C]^\rho+[r(R-L)+C]^\rho\}\\
&\ge c_-(q^\rho+r^\rho)(R+C)^\rho
=c_-(R+C)^\rho.
\end{aligned}
\]

For a rigorous propagation, put \(Q_* =\max(q,r)<1\) and
\(d=(1-Q_*)B+Q_*L>0\). Every child of every \(R\ge B\) is at most
\(R-d\). Induction on the intervals \([B+kd,B+(k+1)d]\), beginning with
the already controlled seed, proves the comparison for all larger radii.
Continuity deals with common endpoints. All power comparisons use only
monotonicity of \(x^\rho\) for positive \(x\); no convexity assumption
is made, so the argument also covers \(\rho<1\).

\hypertarget{explicit-finite-upper-constants}{%
\subsection{Explicit finite upper
constants}\label{explicit-finite-upper-constants}}

For all channels set

\[
Q=\max_j\max(q_j,r_j),\quad
m=\min_j\min(q_j,r_j),\quad
\Delta=\log(1/m),\quad
B_+=\max_j\log\frac1{q_jr_j}+\log2.
\]

Choose an integer \(N\ge0\) with \(Q^{-N}\ge3/m\), and put

\[
V_N=(2^{N+1}-1)\log2+
       (2^{N+1}-N-2)\log(1/Q),
\]

\[
c_+=e^{-\rho}(V_N+1+\Delta+B_+).
\]

These constants are intentionally conservative, and require no numerical
evaluation of an unknown solution. The tangent majorant gives
\(M(1)<2\). At \(R_k=Q^{-k}\), the integral equation gives

\[
M(R_k)\le R_k[1+M(R_{k-1})^2].
\]

Define \(U_0=2\), \(U_k=Q^{-k}(1+U_{k-1}^2)\). Then \(M(R_k)\le U_k\),
\(U_k\ge1\), and \(\log U_k\le k\log(1/Q)+\log2+2\log U_{k-1}\). Summing
this elementary recurrence gives \(\log U_N\le V_N\). Since \(e<3\), the
logarithmic-radius seed \(1\le t\le1+\Delta\) has
\(e^t\le e/m<3/m\le R_N\), and therefore \(\log M(e^t)\le V_N\).

Let \(a_j=\log(1/q_j)\), \(b_j=\log(1/r_j)\), and
\(G(t)=\log M(e^t)+t+B_+\). For \(t\ge1+\Delta\), all child moduli are
at least one, so the positive integral and weight normalization give

\[
\log M(e^t)\le t+\log2+
\max_j\{\log M(e^{t-a_j})+\log M(e^{t-b_j})\}.
\]

The definition of \(B_+\) implies

\[
G(t)\le\max_j\{G(t-a_j)+G(t-b_j)\}.
\]

On the seed, \(G(t)\le V_N+1+\Delta+B_+\le c_+e^{\rho t}\). Each channel
satisfies \(q_j^\rho+r_j^\rho\le1\). Induction in intervals of length
\(\min_j\min(a_j,b_j)>0\) therefore propagates the seed comparison,
giving the explicit upper bound

\[
\boxed{\log M(R)\le c_+R^\rho-\log R-B_+\quad(R\ge e).}
\]

Dividing the two bounds by \(R^\rho\), then taking logarithms of their
positive logarithmic-growth bounds, proves all assertions of Theorem A.1.
The theorem leaves open whether \(\log M(R)/R^\rho\) has a limit.
Lattice resonances or logarithmic oscillations are possible research
directions, not results here.

\hypertarget{joint-holomorphic-dependence-on-parameters}{%
\section{Joint holomorphic dependence on
parameters}\label{joint-holomorphic-dependence-on-parameters}}

Fix the number \(J\) of channels. Allow complex weights \(h_j\), without
normalization, and complex contractions satisfying \(|q_j|,|r_j|<1\).
The same formal recurrence defines coefficients

\[
c_n(h,q,r)\in\mathbb Q[h_1,\ldots,h_J,q_1,r_1,\ldots,q_J,r_J].
\]

Each is homogeneous of degree \(n\) in the weights. If a parameter set
has \(|q_j|,|r_j|\le Q<1\) and \(\sum_j|h_j|\le W\), with \(W>0\),
coefficient induction gives

\[
|c_n(h,q,r)|\le W^n d_n(Q),
\]

where \(f_Q(z)=\sum d_n(Q)z^{2n+1}\) is the established diagonal
function. Consequently, on every disk \(|z|\le R\),

\[
\sum_{n\ge0}|c_n(h,q,r)|R^{2n+1}
\le W^{-1/2}f_Q(\sqrt W R)<\infty.
\]

Every compact subset of the complex parameter domain has such \(Q,W\).
The series therefore converges normally on compact sets of the joint
domain, and its sum is jointly holomorphic in the weights, contractions,
and \(z\). This proves parameter continuity at zero weights, although
the growth order can jump there. Complex or signed parameters do not
inherit the positive-bank growth theorem merely from holomorphic
dependence.

\hypertarget{a-quantitative-compact-bound-for-changing-tap-weights}{%
\section{A quantitative compact bound for changing tap
weights}\label{a-quantitative-compact-bound-for-changing-tap-weights}}

Keep fixed real contractions \(0<q_j,r_j\le Q<1\), and let \(h,g\) be
nonnegative normalized weight vectors. Zero weights are permitted. Let
\(d=\|h-g\|_1\). Then for every \(R\ge0\),

\[
\boxed{\sup_{|z|\le R}|F_h(z)-F_g(z)|
\le\frac d4\{R f_Q'(R)-f_Q(R)\}.}
\]

\textbf{Proof.} Regard \(c_n\) as a polynomial in the real weight
variables with the contractions fixed. Its coefficients are nonnegative.
Induction in the recurrence gives the coefficientwise polynomial bound

\[
0\preceq c_n(h)\preceq d_n(Q)\Bigl(\sum_jh_j\Bigr)^n.
\]

Here \(\preceq\) compares every monomial coefficient. Differentiation
preserves this comparison, so on the normalized nonnegative simplex

\[
0\le\partial_{h_j}c_n\le n d_n(Q).
\]

Put \(v=h-g\). Its positive and negative parts both have sum \(d/2\),
since \(\sum_jv_j=0\). At every point of the segment between \(g\) and
\(h\), all partial derivatives lie in the interval \([0,n d_n(Q)]\).
Hence

\[
|v\cdot\nabla c_n|\le\frac d2 n d_n(Q),\qquad
|c_n(h)-c_n(g)|\le\frac d2 n d_n(Q).
\]

Multiply by \(R^{2n+1}\) and sum. Entire convergence justifies the sum,
and
\[R f_Q'(R)-f_Q(R)=2\sum_{n\ge0}n d_n(Q)R^{2n+1},\]
proving the
claim. The same argument applies to the closed simplex and needs no
positive minimum weight.

\hypertarget{a-vanishing-tap-and-discontinuous-order}{%
\subsection{A vanishing tap and discontinuous
order}\label{a-vanishing-tap-and-discontinuous-order}}

Take the two diagonal channels \((q,r)=(1/4,1/4)\) and \((3/4,3/4)\),
with weights \(1-\varepsilon\) and \(\varepsilon\), respectively. For
\(0\le\varepsilon\le1\), write the resulting function as
\(F_\varepsilon\). The quantitative bound gives

\[
\sup_{|z|\le R}|F_\varepsilon(z)-f_{1/4}(z)|
\le\frac{\varepsilon}{2}\{R f_{3/4}'(R)-f_{3/4}(R)\}.
\]

Thus convergence to \(f_{1/4}\) is uniform on every fixed disk, at an
explicit linear rate in \(\varepsilon\). Yet

\[
\rho(F_0)=\frac12,\qquad
\rho(F_\varepsilon)=\frac{\log2}{\log(4/3)}
\quad(0<\varepsilon\le1).
\]

For the faster active channel the lower constants become
\(L=1/\varepsilon\), \(C=3/\varepsilon\), \(B=1/\varepsilon+8/3\). Its
certified lower-type constant is

\[
c_-(\varepsilon)=\frac{\log2}{(4/\varepsilon+8/3)^{\log2/\log(4/3)}}>0.
\]

It tends to zero as the faster tap vanishes. These are proved growth and
compact-convergence statements. They do not determine an actual
crossover radius or leading asymptotic constant.

